
\chapter{Conclusions} \label{CH7_Conclusions}

\section{Overview}

This dissertation set out to address one of the defining operational challenges of
modern scientific computing: the growing mismatch between the energy demands of
high-performance computing infrastructure and the efficiency of the software systems
used to manage those demands.
Over the preceding six chapters, the work progressed from establishing the scale and
origin of the problem, through a series of increasingly general and flexible
reinforcement learning solutions for node-level power management, to a final broadening
of scope that brought energy-aware resource control to the edge-to-HPC computing
continuum.
Taken together, the contributions form a coherent arc: each chapter addressed a
specific limitation of its predecessor, and the cumulative result is a set of
complementary techniques that collectively demonstrate the viability of
learning-based, adaptive power management across diverse HPC and scientific computing
environments.

The following sections revisit each chapter in turn, articulating what was contributed,
what was demonstrated quantitatively, and how each piece connects to the others.
The chapter closes with a synthesis of cross-cutting themes and a discussion of the
broader implications and future directions opened by this body of work.

\section{Background and Motivation}

Chapter~\ref{CH1_Introduction} established the scale and urgency of the energy problem
in AI and data center infrastructure.
Global data centers consumed approximately 415~TWh of electricity in 2024, a figure
comparable to the entire annual consumption of France and projected to nearly double to
945~TWh by 2030 as AI-specific server capacity expands at roughly 30\% per year.
At the level of a single model, training GPT-3 required 1,287~MWh in a single run,
equivalent to 552 metric tons of CO$_2$ equivalent, and inference costs have made
operational energy the dominant concern at deployment scale: a single ChatGPT query
consumes approximately 7.5 times the energy of a conventional web search.
These figures grounded the case that energy efficiency in data centers is a matter of
global resource allocation, not merely an operational convenience.

Tracing this demand inward to the compute node, the chapter identified the fundamental
source of temporal and spatial unevenness in power consumption: the two-order-of-magnitude
latency gap between processor registers and main memory, the SRAM-DRAM technology
mismatch, NUMA asymmetry in multi-socket configurations, and cache coherence overhead
across cores.
The result is that static power budgets are inherently inefficient, motivating the
need for adaptive, runtime power management that responds to the instantaneous behavior
of the executing workload.
Reinforcement learning was introduced as the appropriate framework for designing such a
controller: an RL agent can learn the non-linear relationship between power cap and
application performance through interaction, without requiring either a hand-coded model
of the processor or prior characterization of the specific workload.

Chapter~\ref{CH2_LiteratureReview} provided the theoretical and technical foundations
for the work that followed.
The chapter surveyed the physical origins of the power wall in the breakdown of
Dennard scaling, the energy-aware scheduling theory that provides offline-optimal
baselines through algorithms such as YDS, the hardware actuators available to software
(principally Intel RAPL for CPU power capping and DVFS for frequency scaling), and the
landscape of control and learning strategies proposed in the prior literature, from
classical PID controllers to model-free deep RL.
A critical observation drawn from the survey is that hardware-centric and
application-specific power management approaches each address only half the problem:
the former optimizes at the circuit level without application awareness, while the
latter requires re-parameterization for every new application or platform.
The gap between these two classes of approaches motivated the application-agnostic,
hardware-agnostic design philosophy that runs through all three RL contributions in
this dissertation.
The chapter also covered the performance metrics used to close the feedback loop in
practice, particularly progress-rate signals derived from lightweight application
instrumentation, and surveyed the edge computing landscape relevant to
Chapter~\ref{CH6_Charon}.

\section{Online Reinforcement Learning for Node-Level Power Management}

Chapter~\ref{CH3_IC2E} presented the first reinforcement learning framework for
dynamic, real-time power capping of HPC compute nodes, establishing the feasibility
of RL-driven power management as an alternative to classical control methods.
The central contribution was a model-based online RL approach in which a Proximal
Policy Optimization (PPO) agent is trained on a lightweight mathematical surrogate
model of the compute node.
The surrogate captures the exponential relationship between the applied power cap and
the application's instantaneous progress toward its scientific objective, derived from
static characterization experiments on the target hardware.
Training on the surrogate rather than on the live system eliminated the risk of
degrading production workloads during the exploration phase of learning, a critical
safety requirement in shared HPC environments.

The controller was deployed through Intel's Running Average Power Limit (RAPL)
interface, which exposes per-domain power capping and energy monitoring on modern
Intel processors.
Evaluation on an Intel Cascade Lake node running the STREAM memory-bandwidth benchmark
demonstrated that the PPO agent achieves a 21\% reduction in energy consumption
relative to unconstrained maximum-power execution, while incurring only a 6.5\%
increase in execution time.
This result was obtained with an average energy consumption of 48.23~kJ over a mean
execution time of 261.19 seconds, representing the operating point that the agent
identified as optimal in the energy-performance trade-off space.
A direct comparison against a classical PI controller, the state-of-the-art method at
the time, confirmed that the RL agent achieves a comparable Pareto-frontier position
with the additional advantage of eliminating the need to manually specify a setpoint
tolerance parameter.
Repeatability experiments across seven distinct hardware configurations on Chameleon
Cloud verified that the energy savings are consistent and not hardware-specific.

The limitations of this approach were equally important to articulate.
The surrogate model must be constructed through static characterization for each new
hardware platform, introducing a one-time overhead that scales poorly when the
controller is to be deployed across a heterogeneous fleet of nodes.
Additionally, the trained policy is tied to the specific benchmark used during model
construction, limiting generalization across the diverse portfolio of applications
that characterize real HPC workloads.
These limitations directly motivated the design choices of the following chapter.

\section{Offline Reinforcement Learning for Application-Agnostic Energy Efficiency}

Chapter~\ref{CH4_HPDC} advanced the RL framework by eliminating both the surrogate
model requirement and the single-application constraint, producing a controller that
is genuinely agnostic to both the hardware platform and the running application.
The key insight was that the data needed to train an effective power management policy
need not come from live interaction with the target system: pre-collected execution
traces, gathered by running a set of benchmark applications at a fixed set of power
caps, contain sufficient information to learn a general-purpose policy through offline
RL.
This offline formulation enabled policy development for new hardware platforms without
requiring access to those platforms during training, a meaningful advantage for
environments where live experimentation on production nodes is impractical.

The controller was trained on a dataset of execution traces collected across multiple
benchmark applications spanning both memory-bandwidth-bound and compute-bound
workloads.
At deployment, the controller observes the same lightweight performance signal used in
Chapter~\ref{CH3_IC2E} (the application's instantaneous progress rate) and selects
the RAPL power cap for the next control interval without any application-specific
configuration.
Validation across twelve benchmark applications demonstrated that the offline RL
controller achieves an average energy saving of 20.34\% with an average performance
degradation of only 7.4\%.
The standard deviation of execution time across repeated runs ranged from 0.47 to
2.91~seconds, and the standard deviation of energy consumption ranged from 0.01 to
1.22~kJ, confirming the consistency and robustness of the learned policy.
The controller also generalized successfully to applications not seen during training,
including workloads with phase-varying behavior, demonstrating that the progress-rate
signal encodes workload-agnostic microarchitectural characteristics that are predictive
of the power-performance relationship across a broad class of scientific applications.

Compared to the online approach of Chapter~\ref{CH3_IC2E}, the offline controller
delivers nearly identical average energy savings (20.34\% versus 21\%) while removing
the per-platform model construction step entirely.
This equivalence in outcome, achieved without live interaction, represents a significant
operational simplification for HPC administrators.
The remaining limitation, however, is that the learned policy encodes a single fixed
energy-performance trade-off, determined implicitly by the reward function used during
training.
In practice, different users and different computational jobs have different tolerance
for performance degradation: a long-running simulation may tolerate a modest slowdown
for substantial energy savings, while a time-sensitive job may require near-maximum
performance regardless of energy cost.
Deploying a separate RL agent for each desired operating point would negate the
generality of the offline training approach.
This limitation defined the research question addressed in the next chapter.

\section{Decay-Driven Multi-Objective Reinforcement Learning for User-Steerable Trade-offs}

Chapter~\ref{CH5_ICS} introduced the Decay-Driven Multi-Objective Reinforcement
Learning (DDMORL) framework, which resolved the single-policy limitation of the
previous two chapters by enabling a single offline-trained controller to cover the
entire energy-performance trade-off surface through a user-specified preference
parameter.
The central technical contribution is the formulation of a preference vector
$\boldsymbol{\omega}$ derived from a user-supplied degradation tolerance
$\delta \in [0,1]$, where $\delta = 0$ corresponds to maximum performance preservation
and $\delta = 1$ corresponds to maximum energy reduction.
This scalar is converted to a preference-aligned control signal via a lightweight
application-specific progress-to-power-cap model $\mathcal{G}_{app}$, which maps the
desired degradation level to the range of RAPL power cap values that are likely to
achieve it.
At deployment, the controller selects actions using a preference-aware scoring function
that jointly maximizes alignment with the user's specified preference and the expected
cumulative reward estimated by the offline-trained Q-network.

The DDMORL framework trains a single Q-network offline on a mixed dataset of eight
benchmarks comprising five STREAM memory-bandwidth variants and three NAS Parallel
Benchmark variants.
No online exploration, no system downtime for retraining, and no enumeration of static
power caps are required at deployment time.
Evaluation on nine benchmarks confirmed three key properties.
First, the preference knob $\delta$ reliably steers the controller across the full
energy-performance trade-off surface: PCAP selection shifts monotonically from the
high end of the action space (approximately 130 to 165~W) for $\delta \leq 0.2$ to the
low end (approximately 89 to 107~W) for $\delta \geq 0.7$, covering all 14 discrete
power-cap levels in the action space $\{89, 95, \ldots, 165\}$~W.
Second, the STREAM benchmarks achieve energy savings of 31 to 39\% with well-separated
Pareto operating points across the full 3 to 95\% degradation range, while the NPB
benchmarks achieve 20 to 31\% savings, with operating-point clustering in the
high-degradation regime reflecting the saturation behavior of compute-bound workloads
at low power caps.
Third, overlaying DDMORL operating points against static fixed-PCAP baselines in the
execution-time versus energy plane shows that the controller's output is
indistinguishable from that of an optimally tuned static power cap, providing empirical
evidence that the preference inference step introduces no measurable runtime or energy
overhead.

The framework also demonstrated out-of-distribution generalization: trained on a dataset
that held NPB-FT out entirely, the controller produced coherent and monotone Pareto
fronts for this unseen benchmark, a result attributed to the PAPI-derived state features
(instructions per cycle, L3 cache miss rate, and stall ratio) that encode
workload-agnostic microarchitectural characteristics predictive of the application's
response to power cap changes.
This result confirms that DDMORL learns a preference-conditioned power management
strategy rather than a workload-specific lookup table, providing a strong foundation
for deployment in the heterogeneous multi-application environments that characterize
production HPC facilities.

Taken together, the three chapters on node-level RL form a coherent progression: from
a model-based online controller requiring per-platform characterization (Chapter~\ref{CH3_IC2E}),
to a model-free offline controller that is application-agnostic but single-policy
(Chapter~\ref{CH4_HPDC}), to a model-free offline controller that is both
application-agnostic and user-steerable across the full trade-off surface
(Chapter~\ref{CH5_ICS}).
Each step removed a real operational barrier to deployment without sacrificing the
efficiency gains demonstrated by its predecessor.

\section{Edge-to-HPC Resource Management for Scientific Workflows}

Chapter~\ref{CH6_Charon} broadened the scope of the dissertation from per-node power
management within a single HPC facility to dynamic resource management across the
edge-to-HPC computing continuum.
The motivation for this extension is that modern large-scale scientific experiments,
particularly those in experimental physics, materials science, and high-energy
astrophysics, generate data at rates that make end-to-end transmission to HPC centers
impractical.
Instead, they rely on an integrated pipeline in which data is generated at remote
sensors or instruments, processed at the edge to reduce volume and extract features,
and then forwarded to HPC systems for full-scale simulation and analysis.
Managing the resources of this edge-HPC continuum under real-time scientific performance
requirements introduces a resource arbitration problem that has no direct analogue in
the single-node power management literature.

The Charon framework, introduced in this chapter, addressed three specific research
questions: how edge resources can be arbitrated dynamically to ensure performance for
experiment-critical applications while remaining flexible for reconfiguration, how
system software can support AI-at-the-edge applications that leverage connected HPC
platforms for model tuning, and how flexible user-agnostic system software policies can
be deployed across diverse hardware setups and use cases.
The framework was validated on a representative scientific use case: a synchrotron
beamline experiment requiring real-time processing of X-ray diffraction images at a
throughput target of 1,000 images per second.
A proportional-derivative (PD) controller was designed to arbitrate edge computational
resources in response to the observed throughput relative to the user-specified target,
and adaptive model quantization was used to adjust the computational cost of the
inference pipeline at runtime.
The experiments confirmed that the PD controller reliably achieved the 1,000 images per
second target, and that model quantization provided a practical mechanism for managing
the trade-off between inference accuracy and computational cost under resource constraints.

The Charon results also identified a significant opportunity for future work: the
PD controller, while effective for the specific use case studied, does not generalize
automatically to new performance criteria or to workflows with complex multi-stage
pipelines.
Replacing the classical controller with a reinforcement learning agent in the control
loop would bring to the edge-to-HPC setting the same model-free adaptivity that
Chapters~\ref{CH3_IC2E} through~\ref{CH5_ICS} demonstrated at the node level, and
the data collection infrastructure developed for the Charon framework provides a
natural starting point for constructing the training environment such an agent would
require.

\section{Cross-Cutting Themes}

Several themes run through all four technical contributions and are worth articulating
explicitly, as they reflect design principles with broad applicability beyond the
specific systems studied here.

\textbf{Application-agnostic feedback signals.}
All three node-level RL controllers rely on the application's instantaneous progress
rate as the primary performance signal.
This signal is defined as the rate at which the application advances toward its
scientific objective, measured through lightweight instrumentation at pre-identified
progress checkpoints in the application code.
The key advantage of this signal is that it is workload-agnostic in the following sense:
while the absolute value of the progress rate depends on the application, its
sensitivity to changes in the power cap is predictable from microarchitectural features
(IPC, cache miss rate, stall ratio) that are available through hardware performance
counters.
This means that a single trained policy can generalize across applications with
different computational characteristics, provided the feature representation captures
the aspects of the microarchitectural state that determine the power-performance
relationship.

\textbf{The offline RL paradigm for systems control.}
A central methodological contribution of this dissertation is the demonstration that
offline RL, trained on pre-collected execution traces, is sufficient for learning
effective power management policies in HPC environments.
This is a non-trivial result: offline RL is known to struggle with distribution shift
when the deployed policy visits states that are underrepresented in the training dataset.
The success of the approach in Chapters~\ref{CH4_HPDC} and~\ref{CH5_ICS} suggests that
the state-action distribution induced by static power cap sweeps is rich enough to
cover the states visited by the optimal policy, at least for the class of workloads
studied.
This finding has practical significance beyond HPC: it implies that production systems
with access to historical telemetry data can develop adaptive control policies without
any live experimentation, which is a prerequisite for deploying RL in safety-critical
or high-availability computing environments.

\textbf{Preference-conditioned control as a deployment interface.}
The DDMORL framework's contribution of a single user-facing scalar parameter $\delta$
as the interface between the operator's performance requirements and the controller's
internal decision-making reflects a broader principle: that the complexity of a
multi-objective optimization problem should be hidden from the user at the deployment
interface.
HPC administrators and scientific computing users are not equipped to reason about
abstract preference weights in a multi-objective reward formulation, but they can
readily specify a maximum acceptable performance degradation in percentage terms.
Bridging this gap through the $\mathcal{G}_{app}$ progress model demonstrates that
domain-specific knowledge about the system can be used to transform a user-facing
specification into a form that is compatible with the internal RL machinery.

\textbf{The continuum from node to edge.}
Chapters~\ref{CH3_IC2E} through~\ref{CH5_ICS} address energy efficiency at the
granularity of a single compute node, while Chapter~\ref{CH6_Charon} addresses
resource management at the scale of a distributed edge-to-HPC workflow.
These two levels of the hierarchy are not independent: the efficiency of the HPC nodes
that receive data from edge devices is a direct input to the end-to-end performance of
the scientific workflow, and the resource decisions made at the edge affect the volume
and timing of data arriving at the HPC tier.
A fully integrated adaptive control system for the edge-to-HPC continuum would need to
coordinate decisions across both levels simultaneously, a challenge that the individual
contributions of this dissertation are well-positioned to address in combination.

\section{Broader Impact and Future Directions}

The quantitative achievements of this dissertation can be stated compactly.
The online RL controller of Chapter~\ref{CH3_IC2E} demonstrated 21\% energy reduction
with 6.5\% performance loss on a single benchmark.
The offline RL controller of Chapter~\ref{CH4_HPDC} generalized this to 20.34\%
average savings with 7.4\% average degradation across twelve benchmarks, without
requiring any live interaction with the target system during training.
The DDMORL framework of Chapter~\ref{CH5_ICS} extended the result further, achieving
31 to 39\% savings on memory-bandwidth-bound workloads and 20 to 31\% on compute-bound
workloads, while placing the control of the energy-performance trade-off directly in
the hands of the user through a single interpretable parameter.
The Charon framework of Chapter~\ref{CH6_Charon} demonstrated real-time resource
arbitration for a scientific edge-to-HPC workflow, meeting a 1,000 images per second
throughput target through adaptive PD control and model quantization.

The most immediate future direction for the node-level work is the extension of DDMORL
to multi-node environments, where PCAP decisions must be coordinated across nodes that
share communication bandwidth.
The state and action spaces in this setting grow combinatorially with the number of
nodes, and the single-node progress signal must be augmented with inter-node
communication statistics to capture the bottlenecks that arise in tightly coupled
parallel applications.
Hierarchical control architectures, in which the node-level DDMORL controller developed
here operates within a higher-level job-scheduler-level policy, represent a natural
path toward end-to-end energy optimization from the scheduling layer down to the
hardware actuator.

For the edge-to-HPC setting, the most promising direction is the replacement of the
PD controller in Charon with a reinforcement learning agent.
The monitoring infrastructure and performance signal definitions developed for Charon
provide the environment interface such an agent would require, and the offline RL
methodology demonstrated in Chapters~\ref{CH4_HPDC} and~\ref{CH5_ICS} could be
applied directly to the execution traces collected by the Charon framework, avoiding
the need for live experimentation on production beamline systems.

A broader research question opened by this work is the applicability of the progress-rate
feedback paradigm to GPU-accelerated workloads and to AI training and inference
pipelines, which represent an increasing fraction of the energy footprint documented in
Chapter~\ref{CH1_Introduction}.
Extending the instrumentation and control framework to heterogeneous CPU-GPU nodes,
where multiple power domains must be controlled jointly and the progress signal must
account for data transfer latency between host and device memory, is a technically
demanding but high-impact problem.
The success of the CPU-only controllers documented here provides both a baseline and a
methodology that can be adapted to this setting.

Finally, the DDMORL framework's demonstration of preference-conditioned control opens
the possibility of integrating energy-performance trade-off management directly into
the job submission interface of HPC facilities.
If each user can specify a performance degradation tolerance at the time of job
submission, the facility scheduler can use the DDMORL controller to operate each node
at the corresponding operating point, aggregating energy savings across the full
workload mix without any manual tuning by the system administrator.
This would bring the efficiency gains demonstrated in controlled single-node experiments
to the scale of production facilities, where the compounded effect of even modest
per-node savings has measurable impact on the facility's total energy expenditure and
carbon footprint.

\section{Closing Remarks}

The work presented in this dissertation demonstrates that reinforcement learning is a
practical and effective framework for adaptive energy management in high-performance
computing, capable of operating without prior knowledge of the hardware platform or
the running application, adaptable to user-specified performance constraints, and
generalizable across diverse workload types.
The progression from model-based online control to model-free offline control to
preference-conditioned multi-objective control reflects a systematic decomposition of
the barriers to deployment in real HPC environments, with each contribution removing
one such barrier while preserving the efficiency gains of its predecessor.
The extension to the edge-to-HPC continuum establishes that the same principles of
feedback-driven adaptive control are applicable at a broader scale, and identifies a
concrete path toward integrating the node-level and continuum-level contributions into
a unified energy-aware resource management framework for scientific computing.
As the energy footprint of data centers and AI infrastructure continues to grow, the
importance of tools that can close the gap between peak hardware capability and
efficient operational practice will only increase, and the techniques developed in this
dissertation represent a meaningful step in that direction.
